\subsection{定义与记号}

设 G 是在局部紧空间 X 上连续地左作用的拓扑群（GT, III, §2, No.~4）；对于 \(s \in G\) 和 \(x \in X\)，记 sx 为 s 对 x 的变换。我们用 \(\boldsymbol{\gamma}_{\mathbf{X}}(s)\) 或 \(\boldsymbol{\gamma}(s)\) 表示由下式定义的从 X 到 X 的同胚：

\[
\gamma (s) x = s x.\tag{1}
\]

有

\[
\boldsymbol {\gamma} (s t) = \boldsymbol {\gamma} (s) \boldsymbol {\gamma} (t).\tag{2}
\]

若 \(f\) 是定义在 \(X\) 上的函数，则 \(\boldsymbol{\gamma}(s)f\) 通过结构传递来定义，即令 \((\boldsymbol{\gamma}(s)f)(\boldsymbol{\gamma}(s)x) = f(x)\)；换言之，

\[
\big (\boldsymbol {\gamma} (s) f \big) (x) = f (s ^ {- 1} x).\tag{3}
\]

若 \(\mu\) 是定义在 X 上的测度，则 \(\gamma(s)\mu\) 也通过结构传递来定义，从而得到

\[
\langle f, \boldsymbol {\gamma} (s) \mu \rangle = \langle \boldsymbol {\gamma} (s ^ {- 1}) f, \mu \rangle \quad \text {   for   } f \in \mathscr {K} (\mathrm{X}).\tag{4}
\]

换言之，

\[
\int_ {\mathrm{X}} f (x) d (\boldsymbol {\gamma} (s) \mu) (x) = \int_ {\mathrm{X}} f (s x) d \mu (x).\tag{5}
\]

若 A 是 \((\gamma(s)\mu)\)-可积集，则 \(s^{-1}A\) 是 \(\mu\)-可积的，并且

\[
\left(\boldsymbol {\gamma} (s) \mu\right) (\mathrm{A}) = \mu \left(s ^ {- 1} \mathrm{A}\right).\tag{6}
\]

测度 \(\boldsymbol{\gamma}(s)\mu\) 也可以定义为 \(\mu\) 在 \(\boldsymbol{\gamma}(s)\) 下的像。有时，与其写作 \(d(\boldsymbol{\gamma}(s)\mu)(x)\)，不如写作 \(d\mu(s^{-1}x)\)；于是公式 (5) 变为：

\[
\int_ {\mathrm{X}} f (x) d \mu (s ^ {- 1} x) = \int_ {\mathrm{X}} f (s x) d \mu (x);
\]

右端的表达式于是可由左端通过“将 \(x\) 替换为 \(sx\)”得到。

定义 1. --- 设 \(\mu\) 是 X 上的一个测度。

\begin{enumerate}
\def\labelenumi{\alph{enumi})}
\item
  称 \(\mu\) 在 G 下不变，如果 \(\pmb{\gamma}(s)\mu = \mu\) 对每个 \(s\in \mathbf{G}\) 都成立。
\item
  称 \(\mu\) 在 \(\mathbf{G}\) 下相对不变，如果 \(\pmb{\gamma}(s)\pmb{\mu}\) 与 \(\pmb{\mu}\) 成比例对每个 \(s\in \mathbf{G}\) 都成立。
\item
  称 \(\mu\) 在 \(\mathbf{G}\) 下拟不变，如果 \(\pmb{\gamma}(s)\pmb{\mu}\) 与 \(\pmb{\mu}\) 等价对每个 \(s\in \mathbf{G}\) 都成立。
\end{enumerate}

评注。--- 1) 假设 \(\mu\) 不变。则 \(|\mu|\)、\(\mathcal{R}(\mu)\)、\(\mathcal{I}(\mu)\) 都是不变的。若 \(\mu\) 为实测度，则 \(\mu^{+}\) 和 \(\mu^{-}\) 也不变。

\begin{enumerate}
\def\labelenumi{\arabic{enumi})}
\setcounter{enumi}{1}
\tightlist
\item
  假设 \(\mu\) 相对不变且非零。对每个 \(s \in G\)，存在唯一的复数 \(\chi(s)\)，使得
\end{enumerate}

\[
\boldsymbol {\gamma} (s) \mu = \chi (s) ^ {- 1} \mu ,\tag{7}
\]

且函数 \(\chi\) 定义在 \(\mathbf{G}\) 上，是 \(\mathbf{G}\) 在 \(\mathbf{C}^*\) 中的表示，称为 \(\mu\) 的乘子。于是公式 (5) 给出

\[
\int_ {\mathrm{X}} f (s x) d \mu (x) = \chi (s) ^ {- 1} \int_ {\mathrm{X}} f (x) d \mu (x),\tag{8}
\]

而公式 (6) 给出

\[
\mu (s \mathrm{A}) = \chi (s) \mu (\mathrm{A}).\tag{9}
\]

按照上述约定，(7) 也可写成

\[
d \mu (s x) = \chi (s) d \mu (x).\tag{10}
\]

\begin{enumerate}
\def\labelenumi{\arabic{enumi})}
\setcounter{enumi}{2}
\tightlist
\item
  由于 \(|\pmb{\gamma}(s)\pmb{\mu}| = \pmb{\gamma}(s)(|\pmb{\mu}|)\)，说 \(\pmb{\mu}\) 拟不变，等价于说 \(|\pmb{\mu}|\) 拟不变。
\end{enumerate}

若 \(\mu\) 拟不变，且 \(\mu'\) 是 X 上与 \(\mu\) 等价的另一测度，则 \(\gamma(s)\mu'\) 与 \(\gamma(s)\mu\) 等价，因而与 \(\mu\) 等价，进而与 \(\mu'\) 等价，所以 \(\mu'\) 拟不变。因此，说 \(\mu\) 在 G 下拟不变，就是说 \(\mu\) 的等价类在 G 下不变。

\(\mu\) 拟不变的充要条件是：X 的局部 \(\mu\)-零测子集的集合在 G 下不变（Ch. V, §5, No.~5, Th. 2）；或者等价地，对 X 的每个 \(\mu\)-零测紧子集 K 及每个 \(s \in \mathbf{G}\)，sK 都是 \(\mu\)-零测的（同上，评注）。

若 \(\mu\) 拟不变，则 \(\mu\) 的支集在 G 下不变。特别地，若 G 在 X 上传递（A, I, §5, No.~5, Def. 6），则该支集或者为空（若 \(\mu = 0\)），或者等于 X（若 \(\mu \neq 0\)）。

引理 1. --- 设 X、Y、Z 是三个拓扑空间，其中 Y 局部紧。设 \((x,y)\mapsto xy\) 是从 \(X\times Y\) 到 Z 的连续映射，并定义映射 \(x\mapsto u_{x}\)，其值在 \(\mathcal{F}(Y;Z)\) 中，且满足关系 \(u_{x}(y)=xy\)。设 f 是 Z 上取值于 R 或某个巴拿赫空间的连续函数，S 是 f 的支集，\(\mu\) 是 Y 上的测度。假设对每个 \(x_{0}\in X\)，存在 X 中 \(x_{0}\) 的一个邻域 V，使得 \(\bigcup_{x\in V}u_{x}^{-1}(S)\) 在 Y 中相对紧。则：

\begin{enumerate}
\def\labelenumi{\alph{enumi})}
\tightlist
\item
  对每个 \(x \in X\)，\(f \circ u_{x}\) 都是 Y 上具有紧支集的连续函数；
\item
  由 a) 定义的映射 \(x \mapsto \int_{Y} f(xy) d\mu(y)\) 在 X 上连续。
\end{enumerate}

断言 a) 是显然的。证明 b)。由于连续性是局部性质，我们可以归结为 \(\bigcup_{x\in X}u_x^{-1}(S)\) 包含在 Y 的某个紧子集 \(Y'\) 中的情形。由于函数 \((x,y)\mapsto f(xy)\) 在 \(X\times Y\) 上连续，\(f\circ u_x\) 一致趋于 \(f\circ u_{x_0}\)（在 \(Y'\) 上，当 \(x\) 趋于 \(x_0\) 时），（GT, X, §3, No.~4, Th. 3），因此 \(\mu(f \circ u_x)\) 趋于 \(\mu(f \circ u_{x_0})\)。引理得证。

现在回到前面的记号。

命题 1. --- 假设 G 局部紧。设 \(\mu\) 是 X 上非零的相对不变测度。则其乘子 \(\chi\) 是 G 上的连续函数。

事实上，取 \(f \in \mathcal{K}(\mathbf{X})\)，令 \(S\) 为 \(f\) 的支集，令 \(s_0\) 为 \(G\) 中一点，并令 \(V\) 为 \(s_0\) 在 \(G\) 中的紧邻域；则集合

\[
\bigcup_ {s \in V} \gamma (s) ^ {- 1} (S) = V ^ {- 1} S
\]

在 X 中是紧的；由引理 1 和公式 (8)，\(\chi(s^{-1})\langle\mu, f\rangle\) 连续依赖于 \(s\)；若选取 \(f\) 使得 \(\langle\mu, f\rangle \neq 0\)，便知 \(\chi\) 连续。

现在设 G 是在局部紧空间 X 上连续地右作用的拓扑群；对于 \(s \in G\) 和 \(x \in X\)，记 xs 为 s 对 x 的变换。我们用 \(\delta_{\mathbf{X}}(s)\) 或 \(\delta(s)\) 表示由下式定义的 X 上的同胚：

\[
\delta (s) x = x s ^ {- 1}.\tag{1'}
\]

有

\[
\boldsymbol {\delta} (s t) = \boldsymbol {\delta} (s) \boldsymbol {\delta} (t).\tag{2'}
\]

通过结构传递，定义 \(\boldsymbol{\delta}(s)\) 在 X 上的函数和测度上的作用：



\[
\big (\boldsymbol {\delta} (s) f \big) (x) = f (x s)\tag{3'}
\]

\[
\langle f, \delta (s) \mu \rangle = \langle \delta (s ^ {- 1}) f, \mu \rangle\tag{4'}
\]

\[
\int_ {\mathrm{X}} f (x) d \bigl (\pmb {\delta} (s) \mu \bigr) (x) = \int_ {\mathrm{X}} f (x s ^ {- 1}) d \mu (x)\tag{5'}
\]

\[
\big (\delta (s) \mu \big) (\mathrm{A}) = \mu (\mathrm{A} s).\tag{6'}
\]

我们约定用 \(d\mu(xs)\) 代替 \(d(\pmb{\delta}(s)\mu)(x)\)，于是 \((5')\) 变为

\[
\int_ {\mathrm{X}} f (x) d \mu (x s) = \int_ {\mathrm{X}} f (x s ^ {- 1}) d \mu (x).
\]

以类似方式定义 X 上在 G 下不变、相对不变和拟不变的测度。若 \(\mu\) 相对不变，则其乘子 \(\chi\) 由下列公式定义：



\[
\boldsymbol {\delta} (s) \mu = \chi (s) \mu\tag{7'}
\]

\[
\int_ {\mathrm{X}} f (x s) d \mu (x) = \chi (s) ^ {- 1} \int_ {\mathrm{X}} f (x) d \mu (x)\tag{8'}
\]

\[
\mu (\mathrm{A} s) = \chi (s) \mu (\mathrm{A})\tag{9'}
\]

\[
d \mu (x s) = \chi (s) d \mu (x).\tag{10'}
\]

若把与 G 对立的群 \(G^{0}\) 看作通过 \((x, s) \mapsto xs\) 作用于 X，则 \(\mu\) 在 \(G^{0}\) 下相对不变，且乘子仍为 \(\chi\)。

最后，设 \(\mathbf{G}\) 是局部紧群。按照公式 \(\pmb{\gamma}(s)x = sx\)、\(\pmb{\delta}(s)x = xs^{-1}\)，它通过左平移和右平移作用于自身。于是

\[
\boldsymbol {\gamma} (s) \boldsymbol {\delta} (t) = \boldsymbol {\delta} (t) \boldsymbol {\gamma} (s).\tag{11}
\]

以上一切在此均适用，因此在 G 上我们有左不变、右不变、相对左不变、相对右不变、左拟不变和右拟不变测度的概念（不过见第 8、9 号）。

映射 \(x \mapsto x^{-1}\) 是从 \(G\) 到 \(G\) 的同胚。对于函数 \(f\)（定义在 \(G\) 上），定义函数 \(\overset{\vee}{f}\)（在 \(G\) 上）：

\[
\stackrel {\vee} {f} (x) = f (x ^ {- 1}).\tag{12}
\]

对于测度 \(\mu\)（定义在 \(\mathbf{G}\) 上），定义测度 \(\breve{\mu}\)：

\[
\stackrel {\vee} {\mu} (f) = \mu (\stackrel {\vee} {f}) \quad \text { for } f \in \mathcal {K} (\mathrm{G}).\tag{13}
\]

换言之，

\[
\int_ {\mathrm{G}} f (x) d \mu^ {\vee} (x) = \int_ {\mathrm{G}} f (x ^ {- 1}) d \mu (x).\tag{14}
\]

若 \(\mathbf{A}\) 是 \(\check{\mu}\)-可积集，则 \(\mathbf{A}^{-1}\) 是 \(\mu\)-可积的，并且

\[
\check {\mu} (\mathrm{A}) = \mu (\mathrm{A} ^ {- 1}).\tag{15}
\]

我们约定用 \(d\mu(x^{-1})\) 代替 \(d\check{\mu}(x)\)，于是 (14) 变为

\[
\int_ {\mathrm{G}} f (x) d \mu (x ^ {- 1}) = \int_ {\mathrm{G}} f (x ^ {- 1}) d \mu (x).
\]

